% TEMPLATE-MILC-v3.tex - plantilla maestra UNICA de las guias MILC v3.
%
% La usan las tres familias de guias, cada una con su driver:
%   · Grados 6-11  -> scripts/build-guias-g11.py
%   · Bebras       -> scripts/build-guias-bebras.py
%   · Semillero    -> scripts/build-guias-semillero.py
%
% Antes habia tres copias casi identicas (~400 lineas iguales, 6 distintas):
% cada mejora habia que aplicarla tres veces, y en la practica no se hacia
% (el motor de recursos vivia solo en dos de ellas). Lo que varia entre
% familias esta parametrizado en 6 placeholders que cada driver llena
% (se nombran aqui SIN su sintaxis real: el builder reemplaza por texto plano,
% tambien dentro de los comentarios, y un valor multilinea rompe el preambulo):
%
%   HEADER_IZQ        encabezado izquierdo de pagina
%   HEADER_DER        encabezado derecho de pagina
%   PORTADA_ETIQUETA  rotulo sobre el numero grande (GRADO/MOMENTO/MODULO)
%   PORTADA_SERIE     linea de serie de la portada
%   PORTADA_ID        identificador de la guia en la portada
%   PORTADA_CIRCULO   el circulito de grado (°), o vacio si no aplica
%
% El resto de claves las documenta content/guias/_SCHEMA.md. El script valida
% que no queden placeholders sin reemplazar antes de escribir el archivo final.

\documentclass[11pt,letterpaper]{article}
\usepackage[margin=2.0cm]{geometry}
\usepackage[table]{xcolor}
\usepackage{fontspec}
\usepackage[spanish]{babel}
\usepackage{tikz}
% Librerías TikZ para los diagramas de las guías (bloque `recursos:`):
%  · babel  → babel en español vuelve activos caracteres como «>», y TikZ los usa
%             en sus opciones (>=stealth, ->). Sin esta librería, cualquier
%             diagrama con flechas revienta con «\language@active@arg> has an
%             extra }». Debe ir ANTES de que se dibuje nada.
%  · positioning        → sintaxis «below left=2pt and 3pt».
%  · decorations.pathreplacing → llaves ({brace}) para señalar rangos.
\usetikzlibrary{babel,positioning,decorations.pathreplacing}
\usepackage{graphicx}
% Circuitos: el trabajo de electrónica se hace hoy con micro:bit y sus sensores
% integrados (MakeCode, sin cableado), así que no se necesitan esquemáticos.
% Si algún día entran montajes con protoboard, basta descomentar esta línea:
% los diagramas .tex/.tikz declarados en `recursos:` ya se incrustan solos.
% \usepackage{circuitikz}
\usepackage[most]{tcolorbox}
\usepackage{tabularx}
\usepackage{array}
\usepackage{fancyhdr}
\usepackage{titlesec}
\usepackage[document]{ragged2e}  % bandera a la derecha en todo el cuerpo
\usepackage{enumitem}
\usepackage{microtype}

% Helvetica Neue no trae la flecha U+2192, asi que cada «→» escrito literalmente
% en el YAML se compilaba a NADA, en silencio (462 flechas en 61 guias: rutas del
% tipo «paleta → tipografia → lockup» salian rotas). Mapeamos el caracter a su
% version matematica, que si tiene glifo. Asi el autor puede seguir escribiendo
% «→» comodamente en el YAML.
\usepackage{newunicodechar}
\newunicodechar{→}{\ensuremath{\rightarrow}}
% Mismo problema con los simbolos de marca y vinetas: el «✓» de «una palomita ✓»
% salia como cuadrito vacio en el PDF de todas las guias que lo usaban.
\usepackage{pifont}
\newunicodechar{✓}{\ding{51}}
\newunicodechar{✗}{\ding{55}}
\newunicodechar{✔}{\ding{52}}
\newunicodechar{•}{\textbullet}
\newunicodechar{≥}{\ensuremath{\geq}}
\newunicodechar{≤}{\ensuremath{\leq}}

\setmainfont{Helvetica Neue}
\setsansfont{Helvetica Neue}
\setlength{\parindent}{0pt}
\setlength{\parskip}{6pt}
% Sin particion de palabras: en 6.º-7.º una palabra cortada por guion al final
% de linea («Comuni-cacion») frena la lectura mucho mas que una linea corta.
\hyphenpenalty=10000
\exhyphenpenalty=10000
\pretolerance=2000
\tolerance=3000
\setlength{\emergencystretch}{3em}
\linespread{1.10}
\renewcommand{\arraystretch}{1.25}
\setlist{leftmargin=5mm,itemsep=1mm,topsep=1mm}
\newcolumntype{Y}{>{\RaggedRight\arraybackslash}X}

\definecolor{milcMagenta}{HTML}{E6007E}
\definecolor{milcVino}{HTML}{5A0038}
\definecolor{milcTurquesa}{HTML}{0093A5}
\definecolor{milcVerde}{HTML}{58A923}
\definecolor{milcMostaza}{HTML}{E5B400}
\definecolor{milcOcre}{HTML}{B86600}
\definecolor{milcNegro}{HTML}{241816}
\definecolor{milcGris}{HTML}{F7F7F4}
\definecolor{milcLinea}{HTML}{E8E2DD}
\definecolor{milcGradePrimary}{HTML}{0D9488}
\definecolor{milcGradeDark}{HTML}{115E59}
\definecolor{milcGradeSoft}{HTML}{CCFBF1}
\definecolor{milcGradeAccent}{HTML}{CFE7FF}
\definecolor{milcGradeText}{HTML}{FFFFFF}
\color{milcNegro}

\pagestyle{fancy}
\fancyhf{}
\lhead{\small Territorio interior · Educación socioemocional}
\rhead{\small Grado 8° · Momento 10 · MILC v3}
\cfoot{\small\thepage}
\renewcommand{\headrulewidth}{0pt}

\newtcolorbox{titlebox}[1]{
  enhanced,colback=#1,colframe=#1,arc=7mm,boxrule=0pt,
  left=7mm,right=7mm,top=3mm,bottom=3mm,
  before skip=8pt,after skip=8pt,
  fontupper=\sffamily\bfseries\Large\color{white}
}

\newtcolorbox{softbox}[3]{
  enhanced,colback=#2,colframe=#1,borderline west={4pt}{0pt}{#1},
  arc=2mm,boxrule=.4pt,left=4mm,right=4mm,top=3mm,bottom=3mm,
  before skip=6pt,after skip=8pt,title={#3},coltitle=white,
  colbacktitle=#1,fonttitle=\sffamily\bfseries,fontupper=\RaggedRight,
  attach boxed title to top left={xshift=3mm,yshift=-2mm},
  boxed title style={arc=2mm, boxrule=0pt}
}

\newtcolorbox{drawbox}[2]{
  enhanced,colback=white,colframe=#1,arc=4mm,boxrule=1pt,
  left=5mm,right=5mm,top=4mm,bottom=4mm,before skip=8pt,after skip=8pt,
  title={#2},coltitle=white,colbacktitle=#1,fonttitle=\sffamily\bfseries,
  fontupper=\RaggedRight,
  attach boxed title to top left={xshift=4mm,yshift=-2mm},
  boxed title style={arc=3mm, boxrule=0pt}
}

\newtcolorbox{infoband}[1]{
  enhanced,colback=#1!8,colframe=#1!50,arc=2mm,boxrule=.5pt,
  left=4mm,right=4mm,top=2mm,bottom=2mm,before skip=4pt,after skip=4pt,
  fontupper=\small\RaggedRight
}

% Puente narrativo entre secciones (contrato editorial Conéctate).
% Da continuidad: "Acabas de X. Ahora vas a Y porque Z."
% Visualmente: banda izquierda en color del grado, texto en itálica pequeña.
\newtcolorbox{puentebox}{
  enhanced,colback=milcGradeSoft,colframe=milcGradeDark,
  borderline west={3pt}{0pt}{milcGradeDark},
  arc=0mm,boxrule=0pt,left=5mm,right=4mm,top=2.5mm,bottom=2.5mm,
  before skip=4pt,after skip=8pt,
  fontupper=\itshape\small\color{milcNegro}\RaggedRight
}
% La flecha va en modo matematico a proposito: Helvetica Neue NO tiene el glifo
% U+2192, asi que \textrightarrow se compilaba a NADA («Missing character: There
% is no → in font Helvetica Neue») y el puente salia sin flecha. En modo
% matematico la flecha viene de la fuente matematica y si se dibuja.
\newcommand{\puente}[1]{%
  \begin{puentebox}%
    \textcolor{milcGradeDark}{$\rightarrow$}\hspace{2mm}#1%
  \end{puentebox}%
}

\newcommand{\linead}[1]{\noindent\color{milcLinea}\rule{#1}{0.5pt}\color{milcNegro}}
\newcommand{\respuesta}[1]{\par\vspace{2mm}\foreach \n in {1,...,#1}{\noindent\color{milcLinea}\rule{\linewidth}{0.5pt}\par\vspace{5mm}}\color{milcNegro}}
\newcommand{\checkbox}{\textcolor{milcGradeDark}{\fbox{\phantom{x}}}\hspace{5pt}}

% ─── Listas aireadas para las guías socioemocionales ────────────────────
% Componer las verificaciones como «\checkbox texto\\» deja el texto que salta
% de línea debajo del cuadrito y apelmaza el bloque. Estos entornos alinean la
% continuación y airean los ítems. Los usa Territorio Interior; cualquier
% familia puede hacerlo.
\newlist{checklista}{itemize}{1}
\setlist[checklista]{
  label=\textcolor{milcGradeDark}{\fbox{\phantom{x}}},
  leftmargin=9mm, labelsep=3.2mm, itemsep=2.4mm,
  topsep=2.5mm, parsep=0pt, partopsep=0pt, before=\RaggedRight
}
\newlist{pasoslista}{enumerate}{1}
\setlist[pasoslista]{
  label=\textcolor{milcGradeDark}{\sffamily\bfseries\arabic*.},
  leftmargin=8mm, labelsep=3mm, itemsep=2.8mm,
  topsep=1mm, parsep=0pt, partopsep=0pt, before=\RaggedRight
}

\newcommand{\phasecard}[4]{
\begin{tcolorbox}[enhanced,colback=#3,colframe=#2,arc=3mm,boxrule=.5pt,height=3.05cm,valign=center,halign=center,left=1.5mm,right=1.5mm,top=2mm,bottom=2mm]
{\sffamily\bfseries\fontsize{22}{24}\selectfont\textcolor{#2}{#1}}\par
{\sffamily\bfseries\small #4}
\end{tcolorbox}
}

% ── Piezas del rediseno 2026 (legibilidad 6.º-11.º) ───────────────────
% Banda de estacion: reemplaza el titlebox macizo. Titulo a la izquierda,
% dato practico (tiempo/modalidad) a la derecha, en una sola linea.
\newcommand{\milcestacion}[2]{%
\begin{tcolorbox}[enhanced,colback=#1,colframe=#1,arc=2mm,boxrule=0pt,
  left=5mm,right=5mm,top=2.6mm,bottom=2.6mm,before skip=11pt,after skip=7pt]
{\sffamily\bfseries\fontsize{15}{18}\selectfont\color{white}#2}
\end{tcolorbox}}

% Tarjeta de paso numerada: sustituye las tablas «Pilar / Aplicacion», que
% eran formato de consulta docente, no de estudiante. El numero va en un
% circulo sobre el borde para que la secuencia se lea de un vistazo.
\newcommand{\milcpaso}[3]{%
\begin{tcolorbox}[enhanced,colback=white,colframe=milcLinea,arc=1.5mm,boxrule=.9pt,
  left=14mm,right=4mm,top=3mm,bottom=3mm,before skip=4pt,after skip=4pt,
  fontupper=\RaggedRight,
  overlay={\node[anchor=north west,circle,fill=#2,text=white,inner sep=0pt,
    minimum size=7.4mm,font=\sffamily\bfseries\small]
    at ([xshift=3.6mm,yshift=-3.2mm]frame.north west) {#1};}]
#3
\end{tcolorbox}}

% Casilla marcable de tamano util con lapiz (la anterior era \fbox{\phantom{x}},
% de alto variable segun el texto que la seguia).
\newcommand{\milcmarca}{\framebox[3.8mm]{\rule{0pt}{3.4mm}}}

% Fila marcable de lista suelta (checks de la estacion 1).
\newcommand{\milccheckrow}[1]{%
\par\vspace{1.8mm}\noindent
\makebox[6.4mm][l]{\raisebox{-.55mm}{\textcolor{milcNegro}{\milcmarca}}}%
\parbox[t]{\dimexpr\linewidth-6.4mm\relax}{\RaggedRight #1}\par}

% Celda marcable dentro de la rubrica (tabla de autoevaluacion). Misma
% sangria francesa, pero pensada para vivir dentro de una columna X.
\newcommand{\milcopcion}[1]{%
\makebox[5.2mm][l]{\raisebox{-.55mm}{\textcolor{milcNegro}{\milcmarca}}}%
\parbox[t]{\dimexpr\linewidth-5.2mm\relax}{\RaggedRight #1}}

% ── Recursos visuales de la guía ──────────────────────────────────────
% Declarados en el bloque `recursos:` del YAML y emitidos por el builder.
% \guiaFigura   → imagen rasterizada o PDF (foto, carta celeste, montaje).
% \guiaDiagrama → fuente TikZ/circuitikz (.tex) incrustada como vector:
%                 es la vía para circuitos y esquemáticos editables.
% #1 = ruta relativa al .tex · #2 = pie (puede ir vacío)
\newcommand{\guiaFigura}[2]{%
\begin{center}
\includegraphics[width=0.88\linewidth,height=0.38\textheight,keepaspectratio]{#1}\\[1.5mm]
{\footnotesize\itshape\textcolor{milcNegro!75}{#2}}
\end{center}
}

\newcommand{\guiaDiagrama}[2]{%
\begin{center}
\input{#1}\\[1.5mm]
{\footnotesize\itshape\textcolor{milcNegro!75}{#2}}
\end{center}
}

\begin{document}
\thispagestyle{empty}

% ── Portada ───────────────────────────────────────────────────────────
% Antes: un rectangulo de color a pagina completa. Se veia bien en pantalla,
% pero en la fotocopiadora del colegio se comia el toner y salia gris sucio.
% Ahora la banda ocupa el tercio superior y el resto va en blanco.
\begin{tikzpicture}[remember picture,overlay]
  \path[fill=milcGradePrimary]
    (current page.north west) rectangle ([yshift=-10.6cm]current page.north east);

  \node[anchor=north west,text=milcGradeText] at ([xshift=2.0cm,yshift=-1.55cm]current page.north west)
    {\sffamily\bfseries\fontsize{12}{14}\selectfont MOMENTO};
  \node[anchor=north west,text=milcGradeText,opacity=.85] at ([xshift=2.0cm,yshift=-2.35cm]current page.north west)
    {\sffamily\bfseries\fontsize{8.5}{10}\selectfont TERRITORIO INTERIOR · SOCIOEMOCIONAL};
  \node[anchor=north east,text=milcGradeText,opacity=.85,align=right] at ([xshift=-2.0cm,yshift=-1.55cm]current page.north east)
    {\sffamily\bfseries\fontsize{9}{12}\selectfont I.E. Sor María Juliana\\Cartago, Valle del Cauca};

  \node[anchor=west,text=milcGradeText] (gradonum) at ([xshift=1.85cm,yshift=-7.1cm]current page.north west)
    {\sffamily\bfseries\fontsize{112}{112}\selectfont 10};
  % El circulito de grado (°) solo aplica a las guias de grado («11°»); en
  % Bebras y Semillero el numero grande es un momento/modulo, no un grado.
  % Se posiciona relativo al nodo `gradonum`, no en coordenadas absolutas.
  

  \node[anchor=west,text=milcGradeText,text width=11.4cm,align=left] at ([xshift=1.5cm,yshift=0.5cm]gradonum.east)
    {
     {\sffamily\bfseries\fontsize{9}{11}\selectfont\MakeUppercase{Territorio interior · Socioemocional}}\\[3mm]
     {\sffamily\bfseries\fontsize{21}{25}\selectfont\setlength{\baselineskip}{27pt}La minga\\[4pt]del curso\\[4pt]entrenar, no prometer}\\[3.5mm]
     {\sffamily\bfseries\fontsize{15}{18}\selectfont Grado 8° · Momento 10}
    };
\end{tikzpicture}

\vspace*{9.4cm}
\vfill

{\sffamily\bfseries\fontsize{8}{10}\selectfont\textcolor{milcGradeDark}{LO QUE TE LLEVAS HOY}}

\begin{tcolorbox}[enhanced,colback=white,colframe=milcNegro,arc=2mm,boxrule=1.6pt,
  left=5mm,right=5mm,top=4mm,bottom=4mm,before skip=3pt,after skip=12pt,
  fontupper=\RaggedRight\fontsize{11}{15.5}\selectfont]
La minga del curso, hecha: una acción colectiva de cuidado elegida por el grupo, con tarea para cada quien, comida compartida, registro de lo que efectivamente ocurrió y —lo que la vuelve entrenamiento— tres revisiones a lo largo de cuatro semanas.
\end{tcolorbox}

\begin{tcolorbox}[enhanced,colback=milcGris,colframe=milcGris,arc=2mm,boxrule=0pt,
  left=5mm,right=5mm,top=3.5mm,bottom=3.5mm,after skip=12pt]
{\sffamily\bfseries\fontsize{8}{10}\selectfont\textcolor{milcNegro!55}{ESTA GUÍA ES DE}}\\[3mm]
\begin{tabularx}{\linewidth}{X p{3.2cm} p{3.2cm}}
\linead{\linewidth} & \linead{3.0cm} & \linead{3.0cm}\\[-1mm]
{\fontsize{8.5}{10}\selectfont\textcolor{milcNegro!55}{Nombre completo}} &
{\fontsize{8.5}{10}\selectfont\textcolor{milcNegro!55}{Grupo}} &
{\fontsize{8.5}{10}\selectfont\textcolor{milcNegro!55}{Fecha}}\\
\end{tabularx}
\end{tcolorbox}

\vfill

\noindent\textcolor{milcLinea}{\rule{\linewidth}{1pt}}\\[2mm]
{\sffamily\bfseries\fontsize{9.5}{12}\selectfont Autor: Dr. Álvaro Cárdenas Orozco}\\[1mm]
{\sffamily\fontsize{8.5}{11}\selectfont\textcolor{milcNegro!55}{Metodología MILC · Apertura ancestral · Acción colectiva de cuidado · Triángulo Dussel-Estoicismo-Floridi}}

\newpage

\section*{Apertura: La minga del curso: la empatía que se entrena, no la que se promete}

\begin{softbox}{milcGradeDark}{milcGradeSoft}{Saber ancestral}
Todo este grado termina donde empezó el país: en la \textbf{minga}. El \textbf{CRIC} ---Consejo Regional Indígena del Cauca--- la nombra como una de sus \textbf{prácticas de unidad}, y hoy la palabra se usa tanto para un trabajo colectivo como para un encuentro donde lo que se trabaja es el pensamiento. \textbf{Recoge las tres cosas que has visto todo el año.} \textbf{Se convoca:} alguien pasa la voz \textbf{uno por uno} ---nadie llega a una minga porque «se enteró»--- y eso es lo contrario de dejar a alguien por fuera. \textbf{Se reparte:} hay tarea para cada quien \textbf{según lo que puede hacer}, los mayores una, los niños otra, el que no puede con el peso lleva el agua o hace la comida; \textbf{y nadie sobra}. \textbf{Se come junto:} el que convoca pone la comida, \textbf{y la olla es una sola}. \emph{Y hay una cuarta cosa que casi nunca se cuenta y es la que hace que esto no sea una jornada bonita:} \textbf{la minga se devuelve}. \emph{El que hoy recibió el trabajo de todos, mañana va a la de otro} ---la cuenta no se lleva en dinero sino en turnos, y por eso la minga no es un favor: \textbf{es un sistema}, y dura porque vuelve---.
\end{softbox}

\begin{softbox}{milcTurquesa}{milcGris}{Saber contemporáneo}
¿La empatía se puede entrenar, o se tiene o no se tiene? \textbf{Se midió.} \textbf{Emily Teding van Berkhout} y \textbf{John Malouff} reunieron \textbf{18 estudios con asignación aleatoria} ---el tipo de estudio más exigente--- con un total de \textbf{1.018 participantes}, y evaluaron programas de \textbf{entrenamiento en empatía}. \textbf{El resultado: los programas funcionan, con un efecto de tamaño medio} (\emph{g} = 0,63, que baja a 0,51 al corregir por posible sesgo de publicación) (Teding van Berkhout y Malouff, 2016). \textbf{Es decir: la empatía \emph{es entrenable}}, igual que un músculo o un idioma. \emph{Pero fíjate en dos detalles del estudio, porque son los que le dan forma a esta guía.} \textbf{Primero: lo que se entrena son \emph{programas}} ---secuencias con ejercicios y práctica repartida en el tiempo---, \textbf{no charlas}. \textbf{Segundo:} en esta clase de investigación, \emph{la evidencia más fuerte es la que se mide en \textbf{conducta observable} y no en lo que la gente dice de sí misma}. \textbf{Traducido, y es el criterio con que se evalúa este momento:} \emph{la empatía no se demuestra diciendo que a uno le importa la gente. Se demuestra en lo que uno hizo, y en si lo siguió haciendo cuatro semanas después.}
\end{softbox}

\begin{softbox}{milcMostaza}{milcGris}{Pregunta-puente}
\textit{En la minga \textbf{se convoca uno por uno, nadie queda sin tarea y la cuenta se lleva en turnos}. \textbf{¿Qué necesita realmente alguien de tu curso ---o de tu colegio--- que ustedes podrían resolver} \ldots \textbf{y que llevan meses viendo sin hacer nada}?}
\end{softbox}

\begin{softbox}{milcVerde}{milcGris}{Saber y saber hacer}
Al terminar podrás: (1) \textbf{explicar} por qué la empatía es entrenable y por qué se mide en conducta y no en intenciones; (2) \textbf{convocar} una acción colectiva a la manera de la minga: uno por uno, con tarea para todos y comida compartida; (3) \textbf{repartir} el trabajo según lo que cada quien puede hacer, sin dejar a nadie sin puesto; (4) \textbf{sostener} la acción con tres revisiones a lo largo de cuatro semanas ---que es lo único que separa un entrenamiento de una jornada bonita---.
\end{softbox}



\small
\begin{tabularx}{\textwidth}{>{\bfseries}p{3.0cm}Y}
\rowcolor{milcGradePrimary!12}Autor & Dr. Álvaro Cárdenas Orozco\\
DBA & Comprende los fenómenos que afectan a su territorio, regula sus emociones ante ellos y acompaña a otros con empatía (transversal 6°--11°, articulado con la política «Escuela, territorio de vida» del MEN, Resolución 006519 de 2025, y la Ley 1620 de 2013)\\
Referentes & Primera ayuda psicológica (OMS, 2011/2012) · Directrices IASC (2007) · USGS y Servicio Geológico Colombiano · Pueblo nasa y la minga (CRIC) · Dussel · Estoicismo · Floridi\\
Producto final & La minga del curso, hecha: una acción colectiva de cuidado elegida por el grupo, con tarea para cada quien, comida compartida, registro de lo que efectivamente ocurrió y —lo que la vuelve entrenamiento— tres revisiones a lo largo de cuatro semanas.\\
\end{tabularx}
\normalsize

\puente{Este es el último momento del grado, y no trae herramienta nueva: \textbf{pone a trabajar juntas las nueve anteriores} ---pertenecer, negarse, mirar, el rumor, la ruta, la marca, el cuerpo, la cerca y el liderazgo---.}

\milcestacion{milcGradeDark}{Ruta MILC: así vamos a aprender}
\begin{tabularx}{\textwidth}{>{\centering\arraybackslash}X>{\centering\arraybackslash}X>{\centering\arraybackslash}X>{\centering\arraybackslash}X}
\phasecard{1}{milcVerde}{milcGris}{Escucha\\[-1mm]\small Observo, escucho y pregunto.} &
\phasecard{2}{milcTurquesa}{milcGris}{Sistematización\\[-1mm]\small Convoco y sostengo.} &
\phasecard{3}{milcMagenta}{milcGris}{Praxis\\[-1mm]\small Construyo y aplico.} &
\phasecard{4}{milcVino}{milcGris}{Evaluación\\[-1mm]\small Reflexión liberadora.}
\end{tabularx}


\puente{Primero, lo que de verdad hace falta. \emph{No lo que quedaría bonito: lo que alguien necesita.}}

\milcestacion{milcVerde}{Estación 1 de 3 · Escucha}

\begin{softbox}{milcVerde}{milcGris}{Escena para escuchar}
\textbf{Actividad 1 · IDENTIFICA --- Lo que hace falta.} \textbf{Sin nombres si se trata de personas.} \textbf{Primera parte, individual.} Escribe \textbf{tres cosas que hacen falta} y que ustedes podrían resolver: en el salón, en el colegio, en la cuadra. \emph{Piensa en cosas concretas ---un espacio que nadie usa, alguien que se está quedando atrás en una materia, un salón sin nada en las paredes, una señora del barrio que necesita ayuda, basura en un lugar que todos usan---.} \textbf{Segunda parte.} Al frente de cada una: \emph{¿a quién le sirve? ¿Cuánto tiempo tomaría? ¿Qué haría falta?} \textbf{Tercera parte, la que conecta todo el grado.} \emph{¿Hay algo de esta lista que tenga que ver con alguien que quedó por fuera este año ---una cerca del momento 8, alguien que come solo, alguien de quien se dijo algo---?} \textbf{Y la última:} ¿alguna vez su curso intentó algo así? \emph{¿Qué pasó y por qué se murió?}
\end{softbox}

{\sffamily\bfseries\fontsize{8}{10}\selectfont\textcolor{milcNegro!55}{MARCA CUANDO LO HAYAS HECHO}}
\milccheckrow{Escribiste tres necesidades concretas, no generalidades.}
\milccheckrow{Respondiste a quién le sirve, cuánto tomaría y qué haría falta.}
\milccheckrow{Miraste si alguna se conecta con alguien que quedó por fuera este año.}

\begin{infoband}{milcVerde}


\textbf{En tu cuaderno (Actividad 1 · IDENTIFICA).} Título: «Actividad 1. Lo que hace falta». \textbf{Formato:} las tres necesidades con sus tres preguntas + la conexión con alguien que quedó por fuera + qué pasó la última vez que intentaron algo. \textbf{Extensión:} media página. \textbf{Sabes que terminaste cuando} tus tres necesidades son \textbf{cosas que se pueden hacer}, no deseos.
\end{infoband}

\puente{Ya lo tienes. Ahora, por qué la empatía se entrena, por qué se mide en conducta y qué hace que una minga no se deshaga.}

\milcestacion{milcTurquesa}{Fase 2 · Sistematización: entender la minga}

\begin{softbox}{milcTurquesa}{milcGris}{Cuatro claves sobre la minga}
Cuatro claves para que esto no termine en una cartelera y una foto.
\end{softbox}

\milcpaso{1}{milcTurquesa}{\textbf{La empatía se entrena, y hay evidencia.} \textbf{18 estudios aleatorizados, 1.018 participantes, efecto de tamaño medio} (Teding van Berkhout y Malouff, 2016). \emph{Esto tiene una consecuencia que vale para todo el año:} \textbf{ponerse en el lugar del otro no es un don que unos tienen y otros no}. \textbf{Es una capacidad que mejora con práctica} ---como pasar un balón o resolver ecuaciones---. \emph{Y de ahí sale la crítica a la manera habitual de enseñar esto:} \textbf{una charla no entrena nada}. \emph{Lo que entrenó en esos estudios fueron \textbf{programas}: ejercicios, repeticiones y práctica repartida en el tiempo.}}
\milcpaso{2}{milcTurquesa}{\textbf{Se mide en conducta, no en intenciones.} En esta clase de investigación, \emph{la evidencia que más pesa es la que se ve en el comportamiento, no la que la gente reporta sobre sí misma}. \textbf{Y tiene todo el sentido:} \emph{cualquiera puede contestar en una encuesta que le importan los demás}. \textbf{Por eso este momento no se evalúa por lo que digan ni por lo que sientan:} \emph{se evalúa por lo que quede hecho, y por si lo siguieron haciendo}. \textbf{Y hay una razón práctica todavía más fuerte:} \emph{a la persona que se beneficia de una minga no le sirve la buena intención de nadie ---le sirve que el trabajo esté hecho---.}}
\milcpaso{3}{milcTurquesa}{\textbf{Nadie sin tarea, y ese es el punto ---no un detalle---.} En una minga se reparte \textbf{según lo que cada quien puede}: hay tareas pesadas y tareas livianas, y las dos hacen falta. \emph{Si en la minga de su curso alguien queda mirando, la actividad acaba de reproducir el problema del momento 9 del grado 7} ---y con el agravante de que ahora es delante de todos---. \textbf{Regla concreta y verificable:} \emph{escriban la lista completa del curso y pongan una tarea al frente de cada nombre}. \textbf{Si sobra un nombre, la minga no está lista.} \emph{Y acuérdense del momento 9 de este grado: el que coordina no puede quedarse con la parte más grande.}}
\milcpaso{4}{milcTurquesa}{\textbf{Cuatro semanas, no un día.} \emph{Casi todas las actividades de este tipo mueren igual:} se hacen, se toma la foto, y a los quince días nadie se acuerda. \textbf{Lo que las convierte en entrenamiento es la repetición espaciada}, y por eso este producto \textbf{exige tres revisiones a lo largo de cuatro semanas}. \emph{En cada revisión se contestan tres preguntas y ya:} \textbf{¿sigue en pie lo que hicimos? ¿Quién no ha cumplido su parte y por qué? ¿Qué hay que ajustar?} \emph{Y hay una razón adicional, muy de esta guía:} \textbf{la minga se devuelve.} \emph{Sostener cuatro semanas es la versión escolar de que la cuenta se lleve en turnos y no en favores.}}

\begin{drawbox}{milcTurquesa}{Cómo se arma una minga que no se deshaga}
\begin{pasoslista}
\item \textbf{Se escoge una necesidad real,} de alguien concreto, no una causa general.
\item \textbf{Se convoca uno por uno.} \emph{Nadie se entera por el grupo: a cada quien se le habla.}
\item \textbf{Se reparte con la lista completa del curso al frente.} \emph{Si sobra un nombre, falta una tarea.}
\item \textbf{Se come junto.} \emph{Suena secundario y no lo es: es lo que hace que sea una minga y no un trabajo asignado.}
\item \textbf{Se ponen las tres fechas de revisión el mismo día,} antes de que nadie se vaya.
\end{pasoslista}
\end{drawbox}

\begin{softbox}{milcMostaza}{milcGris}{Errores comunes y su corrección}
\textbf{Escoger una causa general:} «cuidar el planeta» no se puede repartir; recoger la basura del patio el jueves, sí. \textbf{Convocar por el grupo de chat:} en la minga se avisa uno por uno. \textbf{Dejar gente mirando:} reproduce el momento 9 del grado 7 delante de todos. \textbf{Que el coordinador haga la mitad:} eso es trabajar con público. \textbf{Hacerlo «por» alguien y no «con» alguien:} si es para una persona del curso, esa persona también tiene tarea. \textbf{No poner las fechas de revisión:} sin ellas dura quince días. \textbf{Tomar la foto como si fuera el producto:} el producto es la cuarta semana.
\end{softbox}

\begin{infoband}{milcTurquesa}


\textbf{En tu cuaderno (Actividad 2 · EXPLICA).} Título: «Actividad 2. Entrenamiento, no charla». \textbf{Formato:} explica con tus palabras qué encontró el estudio de Teding van Berkhout y Malouff y por qué eso obliga a que este momento tenga revisiones. \textbf{Extensión:} media página. \textbf{Sabes que terminaste cuando} puedes explicar la diferencia entre \textbf{medir la empatía en conducta y medirla en lo que la gente dice de sí misma}.
\end{infoband}

\puente{Ya saben qué hacer. Es hora de convocarla, repartirla y hacerla ---y de poner desde hoy las fechas de las tres revisiones---.}

\milcestacion{milcMagenta}{Fase 3 · Praxis: hacer la minga}

\begin{softbox}{milcMagenta}{milcGris}{Cómo se convoca y se sostiene una minga}
Esto ya no es un ejercicio de cuaderno. Es la única actividad del año que se evalúa por lo que quedó hecho, y por lo que siga en pie cuatro semanas después.
\end{softbox}

\milcpaso{1}{milcMagenta}{\textbf{Escoger la necesidad (12 min, todo el curso).} Pongan en común las listas y escojan \textbf{una}. \textbf{Tres condiciones:} que \textbf{le sirva a alguien concreto}, que \textbf{se pueda hacer en un día} y que \textbf{necesite a más de diez personas} ---si la pueden hacer tres, no es una minga---. \emph{Y una recomendación que vale más de lo que parece:} \textbf{si alguna de las opciones tiene que ver con alguien que este año quedó por fuera, esa es}.}
\milcpaso{2}{milcMagenta}{\textbf{El reparto (12 min).} Escriban \textbf{la lista completa del curso} y pongan \textbf{una tarea al frente de cada nombre}. \emph{Tareas de todos los tamaños:} conseguir materiales, llevar la comida, pedir el permiso, cargar, limpiar, medir, escribir, tomar el registro, avisar uno por uno. \textbf{Regla de la minga:} \emph{a cada quien según lo que puede}. \textbf{Regla del momento 9:} \emph{el que coordina no lleva la parte más grande.} \textbf{Y la regla que define todo:} \emph{si sobra un nombre, todavía no está lista}.}
\milcpaso{3}{milcMagenta}{\textbf{El día (fuera de clase).} Háganla. \textbf{Y tres cosas que no son adorno:} \emph{convocar uno por uno los días previos} ---no por el grupo---; \textbf{comer juntos} algo, aunque sea mínimo, aportado entre todos; y \textbf{dejar registro} de lo que efectivamente se hizo y de quién hizo qué. \emph{Ese registro es el producto ---no las fotos---.}}
\milcpaso{4}{milcMagenta}{\textbf{Las tres revisiones (4 semanas).} \textbf{Pongan las tres fechas el mismo día de la minga}, antes de que nadie se vaya. \emph{En cada una, cinco minutos y tres preguntas:} \textbf{¿sigue en pie? ¿Quién no cumplió y por qué? ¿Qué se ajusta?} \emph{Y en la última, una cuarta:} \textbf{¿a quién le toca convocar la próxima?} \emph{Porque la minga se devuelve, y eso es lo que la hace durar.}}

\begin{drawbox}{milcMagenta}{Antes de cerrar la minga}
\begin{checklista}
\item La necesidad escogida le sirve a \textbf{alguien concreto} y necesita a más de diez personas.
\item Está escrita \textbf{la lista completa del curso} con una tarea al frente de cada nombre, y \textbf{no sobra ninguno}.
\item Se convocó \textbf{uno por uno}, no por el grupo de chat.
\item Se comió junto, con aporte de todos.
\item Hay \textbf{registro de lo que se hizo y de quién hizo qué} ---no solo fotos---.
\item Las \textbf{tres fechas de revisión} quedaron puestas el mismo día, y en la última se decide a quién le toca convocar la próxima.
\end{checklista}
\end{drawbox}

\begin{softbox}{milcMagenta}{milcGris}{Plantilla de trabajo}
\textbf{La convocatoria, uno por uno.} \emph{«Vamos a hacer \_\_\_\_ el \_\_\_\_. A ti te toca \_\_\_\_. ¿Puedes?»} \\[1mm]
\textbf{El reparto.} \emph{Lista completa del curso · una tarea por nombre · ninguno vacío.} \\[1mm]
\textbf{Las revisiones.} \emph{Fechas: \_\_\_\_ · \_\_\_\_ · \_\_\_\_. Preguntas: «¿sigue en pie? · ¿quién no cumplió y por qué? · ¿qué se ajusta?»} ---y en la última: \emph{«¿a quién le toca la próxima?»}---.
\end{softbox}

\begin{infoband}{milcMagenta}


\textbf{En tu cuaderno (Actividad 3 · CREA).} Título: «Actividad 3. La minga del curso». \textbf{Formato:} la necesidad escogida y por qué + la lista completa con la tarea de cada quien + el registro del día + las tres fechas y lo que se respondió en cada revisión. \textbf{Extensión:} una página, más el registro. \textbf{Sabes que terminaste cuando} \textbf{pasó la tercera revisión} ---no el día de la minga---.
\end{infoband}

\puente{El producto no es un plan: \textbf{es la minga hecha}. \emph{Y su prueba no es el día, sino la cuarta semana.}}

\milcestacion{milcMostaza}{Producto de la sesión}
\textbf{La minga del curso: acción colectiva con tarea para todos, comida compartida y tres revisiones en cuatro semanas}.
\begin{softbox}{milcMostaza}{milcGris}{Criterios de calidad}
(1) La necesidad es concreta, beneficia a alguien identificable y requiere a todo el curso. (2) La lista completa del curso tiene una tarea asignada por nombre, sin nadie sin puesto. (3) La convocatoria se hizo persona por persona, y hubo comida compartida con aporte de todos. (4) Existe registro de lo efectivamente hecho y de quién hizo qué, más allá de las fotos. (5) Las tres revisiones se realizaron en cuatro semanas, y en la última se designó quién convoca la siguiente.
\end{softbox}

\puente{Antes de cerrar, revisen lo único que decide si esto sirvió: \emph{¿queda alguien del curso sin tarea?} Si queda, la minga reprodujo el problema del momento 9 del grado 7.}

\milcestacion{milcVino}{Evaluación liberadora · cinco dimensiones}

\begin{softbox}{milcVino}{milcGris}{Sentido de la evaluación}
Evaluar no es solo revisar una nota. Es reconocer qué aprendí, qué cambió en mí, qué emociones manejé, cómo me relaciono con la ciudad y la comunidad, y qué le dejaré a quien viene detrás.
\end{softbox}

\textbf{Mi reflexión liberadora en cinco dimensiones.} Completa cada frase iniciada:

\small
\begin{tabularx}{\textwidth}{>{\bfseries\RaggedRight\arraybackslash}p{3.4cm}Y>{\RaggedRight\arraybackslash}p{4.6cm}}
\rowcolor{milcVino}\color{white}Dimensión & \color{white}\bfseries Mi respuesta (frame starter) & \color{white}\bfseries Algo que descubrí\\
1. Personal & Lo que más cambió en mí fue \linead{6.5cm} & \linead{4.2cm}\\
2. Emocional & La emoción más fuerte fue \linead{2.5cm} y la regulé \linead{3cm} & \linead{4.2cm}\\
3. Ciudadana & En democracia esto importa porque \linead{6cm} & \linead{4.2cm}\\
4. Local (Cartago) & En mi barrio sería útil para \linead{6.5cm} & \linead{4.2cm}\\
5. Intergeneracional & A mi abuela/abuelo le diría \linead{6.5cm} & \linead{4.2cm}\\
\end{tabularx}
\normalsize

\begin{infoband}{milcVino}
\textbf{Para la próxima sustentación.} Vuelve a esta tabla en seis meses. Verás que las respuestas cambian — no porque lo aprendido cambie, sino porque tú cambias. La evaluación liberadora es una conversación contigo a través del tiempo.
\end{infoband}


\puente{Cerramos el grado con tres voces: el que reclama un derecho, la abeja y el enjambre, y el cuidado de lo que se comparte.}

\milcestacion{milcGradeDark}{Triángulo de pensamiento · cierre filosófico}

\begin{softbox}{milcVino}{milcGris}{Dussel · lente del nosotros}
\textit{«El otro se revela realmente como otro\ldots como el pobre, el oprimido; el que, a la vera del camino, fuera del sistema, muestra su rostro sufriente y sin embargo desafiante: «¡Tengo hambre!, ¡tengo derecho a comer!»»}

{\footnotesize\itshape\color{milcNegro!70}--- Enrique Dussel · Filosofía de la liberación (1977)}

Con esta cita se cierra el grado, y conviene leerla con cuidado por lo que \textbf{no} dice. \emph{El otro no pide caridad: \textbf{reclama un derecho}.} \textbf{Y eso cambia por completo el sentido de una minga:} \emph{no es un acto de bondad de un curso generoso hacia alguien necesitado} ---esa versión, la de la obra de caridad, humilla a quien recibe y le da un aire de superioridad a quien da---. \textbf{Una minga es otra cosa:} \emph{es gente que reconoce que le corresponde}. \textbf{Por eso, si la minga es para alguien del curso, esa persona también tiene tarea:} \emph{no es el objeto de la ayuda, es uno más de los que trabajan.}

\textbf{Nuestra minga, ¿la estamos haciendo «por» alguien o «con» alguien?}
\end{softbox}
\linead{\linewidth}\\[1mm]
\linead{\linewidth}

\begin{softbox}{milcMostaza}{milcGris}{Estoicismo · Marco Aurelio · Meditaciones VI, 54 (c. 175 d.C.) · cuidado interior}
\textit{«Lo que no beneficia al enjambre, tampoco beneficia a la abeja.»}

{\footnotesize\itshape\color{milcNegro!70}--- Marco Aurelio · Meditaciones VI, 54 (c. 175 d.C.)}

Es la frase que resume el grado entero. \emph{Cada momento mostró una versión del mismo error:} \textbf{creer que a uno le puede ir bien mientras al grupo le va mal} ---quedarse en el grupo que cobra peaje, callar cuando presionan, mirar sin hacer nada, repetir un rumor, defender una cerca porque «a mí me va bien con los míos»---. \textbf{Marco Aurelio lo desarma en once palabras.} \emph{Y la minga es la demostración práctica:} \textbf{es imposible hacerla solo}, y al final el que puso el trabajo se lleva algo que no habría podido conseguir por su cuenta ---incluido el derecho a que mañana vayan a la suya---.

\textbf{De lo que hicimos hoy, ¿qué me llevo yo que no habría podido conseguir solo?}
\end{softbox}
\linead{\linewidth}\\[1mm]
\linead{\linewidth}

\begin{softbox}{milcTurquesa}{milcGris}{Floridi · lente de la infoesfera}
\textit{«El florecimiento de las entidades informacionales y de toda la infoesfera debe promoverse preservando, cultivando y enriqueciendo sus propiedades.»}

{\footnotesize\itshape\color{milcNegro!70}--- Luciano Floridi · La ética de la información (2013; trad.)}

\textbf{Preservar, cultivar, enriquecer:} tres verbos que describen exactamente lo que hace una minga, \textbf{y que explican por qué esta guía exige cuatro semanas}. \emph{Preservar es que lo hecho siga en pie; cultivar es volver a revisarlo; enriquecer es que la próxima la convoque otro.} \textbf{Y hay una advertencia sobre lo digital que cierra el año:} \emph{la parte más fácil de una minga hoy es publicarla} ---la foto, la historia, el «qué chévere lo que hicimos»---. \textbf{Publicar no es ninguno de los tres verbos.} \emph{Si el registro que queda es más grande que el trabajo que se hizo, lo que hubo fue una foto con actividad de fondo.}

\textbf{¿Qué es más grande en nuestra minga: lo que quedó hecho o lo que se publicó?}
\end{softbox}
\linead{\linewidth}\\[1mm]
\linead{\linewidth}

\begin{infoband}{milcNegro}
\textbf{Nota para el docente.} Las tres son citas textuales y llevan su referencia debajo. Puedes remitir a la obra en clase.
\end{infoband}

\milcestacion{milcVino}{Mi autoevaluación · marca una casilla por fila}

{\small
\setlength{\extrarowheight}{2.2pt}
\begin{tabularx}{\textwidth}{>{\RaggedRight\arraybackslash\bfseries}p{3.0cm}YYY}
\rowcolor{milcVino}\color{white}\bfseries Criterio & \color{white}\bfseries Lo logré & \color{white}\bfseries En proceso & \color{white}\bfseries Necesito apoyo\\[.6mm]
Comprensión del tema & \milcopcion{Explico las ideas con mis palabras.} & \milcopcion{Reconozco las ideas, falta precisión.} & \milcopcion{Necesito repasar lo básico.}\\[.8mm]
\rowcolor{milcGris}Calidad de la minga & \milcopcion{La minga se hizo, todos tuvieron tarea y sobrevivió a las tres revisiones.} & \milcopcion{Se hizo el día de la actividad, pero nadie volvió a revisar nada.} & \milcopcion{Quedó en una lista de buenas intenciones.}\\[.8mm]
Aplicación y producto & \milcopcion{Mi producto cumple los criterios.} & \milcopcion{Está armado, con ajustes pendientes.} & \milcopcion{Aún está en construcción.}\\[.8mm]
\rowcolor{milcGris}Reflexión 5 dimensiones & \milcopcion{Respondí cada una con honestidad.} & \milcopcion{Respondí algunas con detalle.} & \milcopcion{Mis respuestas son muy generales.}\\[.8mm]
Triángulo de pensamiento & \milcopcion{Conecté las 3 voces con mi trabajo.} & \milcopcion{Conecté al menos una.} & \milcopcion{No las relaciono con mi proceso.}\\[.8mm]
\end{tabularx}}

\vspace{3mm}
\textbf{Hoy entendí que:}\\[1mm]
\linead{\linewidth}\\[1mm]
\linead{\linewidth}

\textbf{Mi compromiso para la próxima sesión es:}\\[1mm]
\linead{\linewidth}\\[1mm]
\linead{\linewidth}

\begin{softbox}{milcMostaza}{milcGris}{Cita de despedida · Marco Aurelio}
\textit{``El obstáculo en el camino se vuelve el camino. Lo que estorba al camino se vuelve camino.''}\\[1mm]
Cada error de hoy, bien observado, se convierte en pieza del camino que pisarás mañana.
\end{softbox}

\small
\begin{tabularx}{\textwidth}{>{\bfseries}p{3.6cm}Y}
\rowcolor{milcGradeDark!12}Nombre completo: & \linead{10cm}\\
Fecha: & \linead{4cm} \quad Curso/grupo: \linead{4cm}\\
Mi firma: & \linead{9cm}\\
\end{tabularx}
\normalsize

\begin{infoband}{milcGradeDark}
\textbf{Para las próximas cuatro semanas.} \textbf{Primero, las tres revisiones} ---en las fechas que ustedes pusieron, con las tres preguntas---. \emph{Y guarden el dato más honesto de todos: cuántos siguieron cumpliendo su parte en la tercera.} \textbf{Ese número, y no el día de la minga, es lo que dice si el curso aprendió algo este año.} \textbf{Segundo, pregunta en tu casa por la minga que sí vieron:} averigua con alguien mayor \textbf{si alguna vez participó en un trabajo entre todos} ---levantar una casa, arreglar un camino, una escuela, un acueducto, la vía de la vereda---: \textbf{quién convocaba, cómo se repartía, quién ponía la comida y qué pasó con lo que hicieron}. \textbf{Trae:} el número de la tercera revisión y lo que te contaron. \emph{Vas a encontrar el cierre perfecto para este grado: casi siempre lo que se construyó así todavía está en pie, y casi siempre quien lo cuenta dice el nombre de la gente que estuvo ---no el suyo---.}
\end{infoband}

\end{document}
